\documentclass[aps, reprint, twocolumn, superscriptaddress, floatfix]{revtex4-2}

% --- Writing Packages (Moved from style.sty) ---
\usepackage{amsmath, amssymb}
\let\Bbbk\relax % Avoid conflict with newtxmath
\usepackage{newtxmath, newtxtext} % Load after amsmath
\usepackage{graphicx}
\usepackage{tikz}
\usepackage{tikz-3dplot}
\usetikzlibrary{arrows.meta, angles, quotes, topaths, shadings, decorations.markings}
\usepackage{booktabs}
\usepackage{tikzorbital}
\usepackage{float}
\usepackage{hyperref}
\hypersetup{
    colorlinks=true,
    linkcolor=blue,
    citecolor=blue,
    urlcolor=blue
}

\usepackage{style} % Style definitions only

\begin{document}

\preprint{PHY3109. COMPUTATIONAL PHYSICS}

\title{A Brief Introduction to Computational Physics}

\author{Jeongbin Jo}
\email{jeongbin033@yonsei.ac.kr}
\affiliation{Department of Physics, Yonsei University, Seoul, 03722, Republic of Korea}

\date{\today}

\begin{abstract}
This report provides a comprehensive overview of fundamental topics covered in the PHY3109 Computational Physics course. Computational physics sits at the intersection of physics, computer science, and applied mathematics, offering numerical methods to solve complex physical systems where analytical solutions are intractable. We explore core numerical techniques including numerical integration (Trapezoidal, Simpson's, and Monte Carlo methods), statistical data analysis, curve fitting, and the propagation of uncertainty. Furthermore, we summarize the numerical solutions to ordinary differential equations (ODEs) using Euler and Runge-Kutta algorithms, concluding with practical applications to physical scenarios such as electric fields and the analysis of signals via Fourier transforms.
\end{abstract}

\maketitle
\thispagestyle{fancytitle}
% --- Watermark (Seal) ---
\addwatermark

\onecolumngrid  % 여기서부터 1단 시작
\vspace{-1cm}
\tableofcontents
\vspace{1cm}    % 목차와 본문 사이 간격 조정 (선택 사항)
\twocolumngrid  % 여기서부터 다시 2단 시작

\section{Introduction}

Physics traditionally relies on theoretical derivation and experimental observation. However, as physical models encompass realistic boundary conditions, external non-linear perturbations, and many-body interactions, their mathematical representations rapidly elude exact analytical solutions. Computational physics bridges this gap between theoretical formulation and empirical data by mapping continuous differential and integral equations into discrete numerical algorithms that are solvable computationally. 

This report provides a comprehensive foundation for the mathematical instruments employed in the PHY3109 Computational Physics course, formulated using Python. From rudimentary calculus to probabilistic integration, multidimensional parameter optimization, and non-linear dynamics, the subsequent sections unpack the mathematical rigor and algorithmic structures of these foundational computational methods. Note that the standard deviation notation $\sigma$ will be utilized distinctly across the statistical and optimization derivations according to empirical physical conventions.

\section{Numerical Integration Algorithms}

Evaluating the definite integral $I = \int_{a}^{b} f(x) \, dx$ is a persistent requirement. When analytical antiderivatives do not exist or are prohibitively complex, numerical quadratures and probabilistic methods are deployed.

\subsection{Newton-Cotes Quadrature and Simpson's Rule}
Deterministic grid-based algorithms approximate the integration domain $[a, b]$ into evenly spaced sub-intervals defined by a step size $h = \frac{b-a}{N}$. While the Trapezoidal rule fits linear polynomials between points, imparting a truncation error of $O(h^2)$, Simpson's $\frac{1}{3}$ rule approximates the local curvature using second-order polynomials (parabolas). By utilizing the Taylor series expansion at the midpoint $x_1$ between $[x_0, x_2]$ (where $h = x_1 - x_0 = x_2 - x_1$), integrating the quadratic interpolation yields:
\begin{equation}
    \int_{x_0}^{x_2} f(x) \, dx \approx \frac{h}{3} \left[ f(x_0) + 4f(x_1) + f(x_2) \right]
\end{equation}
When generalized across $N$ (even) contiguous intervals, the composite Simpson's $\frac{1}{3}$ formula is derived:
\begin{equation}
\begin{aligned}
    \int_{a}^{b} f(x) \, dx &\approx \frac{h}{3} \Bigg[ f(x_0) + f(x_N) \\
    &\quad + 4\sum_{i=1,3,5}^{N-1} f(x_i) + 2\sum_{j=2,4,6}^{N-2} f(x_j) \Bigg]
\end{aligned}
\end{equation}
The local truncation error scales analogously to the fourth derivative of the function, yielding $O(h^5)$ locally and $O(h^4)$ globally, establishing it as the standard for 1D deterministic integration.

\subsection{Monte Carlo Probabilistic Integration}
Grid-based geometries scale exponentially in computational cost with dimensionality (the "curse of dimensionality"). Monte Carlo methods circumvent this by treating the integral as an Expected Value computation. By sampling $N$ statistically independent random coordinates uniformly distributed within a high-dimensional volume $V$, the integral translates to the geometric mean formulation:
\begin{equation}
    I = \int_V f(\mathbf{x}) \, d\mathbf{x} \approx \frac{V}{N} \sum_{i=1}^{N} f(\mathbf{x}_i) = V \langle f \rangle
\end{equation}
The uncertainty of this approximation is not bounded by grid granularity but governed by the Central Limit Theorem. The variance of the approximation scales with the variance of the function itself ($\sigma_f^2$):
\begin{equation}
    \sigma_I \approx V \frac{\sigma_f}{\sqrt{N}}
\end{equation}
Thus, the precision converges as $1/\sqrt{N}$, rendering Monte Carlo unconditionally superior for systems involving massive degrees of freedom, such as statistical mechanics ensembles.

\section{Statistical Data Analysis and Optimization Framework}

Real-world physical measurements intrinsically contain noise. Isolating deterministic signals from stochastic deviations requires rigorous statistical estimators and optimization algorithms.

\subsection{Descriptive Statistics: Mean and Variance}
Before initiating optimization, the intrinsic empirical distribution must be established. For a recorded dataset comprising $N$ individual measurements $\{x_1, x_2, \dots, x_N\}$, the primary central tendency estimator is the arithmetic sample mean $\bar{x}$:
\begin{equation}
    \bar{x} = \frac{1}{N} \sum_{i=1}^{N} x_i
\end{equation}
To quantify the stochastic spread of the data, variance and standard deviation are calculated. Maintaining the universally adopted $\sigma$ notation, the unbiased sample standard deviation $\sigma$ incorporates Bessel's correction $(N-1)$ to compensate for the loss of a degree of freedom when the population mean is approximated by the sample mean:
\begin{equation}
    \sigma^2 = \frac{1}{N-1} \sum_{i=1}^{N} (x_i - \bar{x})^2, \quad \sigma = \sqrt{\sigma^2}
\end{equation}

\subsection{Least Squares Method ($\chi^2$ Minimization)}
To construct a theoretical model $f(x; \vec{\theta})$ that maximally matches the empirical data, one formulates an objective function representing the square of the deviations. If each data point $y_i$ encompasses its own statistical error $\sigma_i$, the deviations are weighted proportionally, defining the rigorous chi-squared statistic $\chi^2$:
\begin{equation}
    \chi^2(\vec{\theta}) = \sum_{i=1}^{N} \left[ \frac{y_i - f(x_i; \vec{\theta})}{\sigma_i} \right]^2
\end{equation}
The aim of regression modeling is strictly to discover the parameter vector $\vec{\theta}$ that minimizes this multidimensional $\chi^2$ landscape. 

\subsection{Gradient Descent Optimization}
While analytical minimization ($\nabla \chi^2 = 0$) exists for purely linear polynomials, complex physical non-linear models depend on numeric iteration. The \textbf{Gradient Descent} algorithm navigates the highly non-linear cost function $J(\vec{\theta}) \equiv \chi^2(\vec{\theta})$. The parameters evolve continuously against the direction of the local gradient manifold:
\begin{equation}
    \theta_j^{(n+1)} = \theta_j^{(n)} - \alpha \frac{\partial J(\vec{\theta}^{(n)})}{\partial \theta_j}
\end{equation}
Here, $\alpha$ is the systematically tuned \textit{learning rate}. If $\alpha$ is too steep, the parameter trajectories will predictably diverge across the minimum basin; conversely, an infinitesimally small $\alpha$ restricts convergence within computational timeframes.

\subsection{Error Analysis and Propagation of Uncertainty}
A measured dependent physical phenomenon often convolves multiple independent measurements $f = f(x_1, x_2, \dots, x_M)$. Errors inherent in individual instruments propagate through equations, necessitating rigorous boundary tracking. By utilizing a first-order multivariate Taylor expansion around the expected values, the generalized combined standard deviation ($\sigma_f$) propagating through the operational architecture can be derived. Considering the completely general case, inclusive of variable interdependence (correlation), the full covariance expansion yields:
\begin{equation}
    \sigma_f^2 = \sum_{i=1}^{M} \left( \frac{\partial f}{\partial x_i} \right)^2 \sigma_{x_i}^2 + 2 \sum_{i<j} \left( \frac{\partial f}{\partial x_i} \right) \left( \frac{\partial f}{\partial x_j} \right) \text{cov}(x_i, x_j)
\end{equation}
In this formulation, the partial derivative $\frac{\partial f}{\partial x_i}$ acts as a "sensitivity coefficient," dictating how strongly the measurement error $\sigma_{x_i}$ amplifies within the final composite outcome. 

In the general computational environment, variable orthogonality (zero covariance) cannot always be guaranteed. For instance, variables interconnected by systematic environmental drift will exhibit non-zero covariance ($\text{cov}(A, B) = \sigma_{AB}$). Using the fully expanded Taylor matrix, Table \ref{tab:error_prop} specifies the generalized variance and standard deviation propagation rules for fundamental arithmetic and continuous physics functions.

\begin{table*}[htp]
    \centering
    \caption{Comprehensive Error Propagation Framework (including covariance $\sigma_{AB}$)}
    \label{tab:error_prop}
    \renewcommand{\arraystretch}{2.2}
    \resizebox{\textwidth}{!}{
    \begin{tabular}{@{}lllc p{5cm}@{}}
        \toprule
        \textbf{Function} $f$ & \textbf{Variance} $\sigma_f^2$ & \textbf{Standard Deviation} $\sigma_f$ \\
        \midrule
        $f = aA$ & $a^2 \sigma_A^2$ & $|a|\sigma_A$ \\
        $f = aA \pm bB$ & $a^2\sigma_A^2 + b^2\sigma_B^2 \pm 2ab\sigma_{AB}$ & $\sqrt{a^2\sigma_A^2 + b^2\sigma_B^2 \pm 2ab\sigma_{AB}}$ \\
        $f = AB$ & $f^2 \left[ \left(\frac{\sigma_A}{A}\right)^2 + \left(\frac{\sigma_B}{B}\right)^2 + 2\frac{\sigma_{AB}}{AB} \right]$ & $|f| \sqrt{\left(\frac{\sigma_A}{A}\right)^2 + \left(\frac{\sigma_B}{B}\right)^2 + 2\frac{\sigma_{AB}}{AB}}$ \\
        $f = \frac{A}{B}$ & $f^2 \left[ \left(\frac{\sigma_A}{A}\right)^2 + \left(\frac{\sigma_B}{B}\right)^2 - 2\frac{\sigma_{AB}}{AB} \right]$ & $|f| \sqrt{\left(\frac{\sigma_A}{A}\right)^2 + \left(\frac{\sigma_B}{B}\right)^2 - 2\frac{\sigma_{AB}}{AB}}$ \\
        $f = aA^b$ & $\left(\frac{fb\sigma_A}{A}\right)^2$ & $\left| \frac{fb\sigma_A}{A} \right|$ \\
        $f = a \ln(bA)$ & $\left(a\frac{\sigma_A}{A}\right)^2$ & $\left| a\frac{\sigma_A}{A} \right|$ \\
        $f = a \log_{10}(bA)$ & $\left(a\frac{\sigma_A}{A \ln(10)}\right)^2$ & $\left| a\frac{\sigma_A}{A \ln(10)} \right|$ \\
        $f = a e^{bA}$ & $f^2(b\sigma_A)^2$ & $|f| \, |b\sigma_A|$ \\
        $f = a^{bA}$ & $f^2(b\ln(a)\sigma_A)^2$ & $|f| \, |b\ln(a)\sigma_A|$ \\
        $f = a \sin(bA)$ & $[ab\cos(bA)\sigma_A]^2$ & $|ab\cos(bA)\sigma_A|$ \\
        $f = a \cos(bA)$ & $[ab\sin(bA)\sigma_A]^2$ & $|ab\sin(bA)\sigma_A|$ \\
        $f = A^B$ & $f^2 \left[ \left(\frac{B}{A}\sigma_A\right)^2 + (\ln(A)\sigma_B)^2 + 2\frac{B\ln(A)}{A}\sigma_{AB} \right]$ & $|f| \sqrt{\left(\frac{B}{A}\sigma_A\right)^2 + (\ln(A)\sigma_B)^2 + 2\frac{B\ln(A)}{A}\sigma_{AB}}$ \\
        \bottomrule
    \end{tabular}
    }
\end{table*}

This systematic propagation fundamentally restricts the precision limit of composite calculations; for instance, analyzing free-fall bounds inherently conflates the timing standard deviation ($\sigma_t$) squared against the dimensional metric deviation ($\sigma_y$), directly modulating the gravitational constant certainty $\sigma_g$.

\section{Numerical Solutions to Differential Equations}

The continuous physical world, codified by Newton's and Maxwell's equations, operates via differential mathematics. Simulation translates these continuous Initial Value Problems (IVPs), $\frac{dy}{dt} = f(t, y)$, into discrete temporal jumps $\Delta t$.

\subsection{From Euler to Runge-Kutta Formulations}
The classical Euler method is structurally derived from a truncated first-order Taylor expansion:
\begin{equation}
    y(t_{n+1}) = y(t_n) + \Delta t \cdot y'(t_n) + O(\Delta t^2)
\end{equation}
This linear extrapolation generates a disastrous systemic accumulation of truncation error corresponding globally to $O(\Delta t)$. To mitigate this, the \textit{Modified Euler Method} (equivalent to a 2nd-order Runge-Kutta or Midpoint Method) evaluates the derivative at the half-step to construct a more accurate predictor-corrector paradigm:
\begin{equation}
\begin{aligned}
    k_1 &= f(t_n, y_n) \\
    k_2 &= f\left(t_n + \frac{\Delta t}{2}, y_n + \frac{\Delta t}{2} k_1 \right) \\
    y_{n+1} &= y_n + \Delta t \cdot k_2 
\end{aligned}
\end{equation}

By extending this logic, the definitive solver in computational mechanics becomes the 4th-order Runge-Kutta (RK4). By rigorously cross-referencing slopes across four distinct spatial-temporal nodes uniformly distributed across the integration step, RK4 yields a highly stable numerical architecture: 
\begin{equation}
\begin{aligned}
    k_1 &= f(t_n, y_n) \\
    k_2 &= f\left(t_n + \frac{\Delta t}{2}, y_n + \frac{\Delta t}{2} k_1\right) \\
    k_3 &= f\left(t_n + \frac{\Delta t}{2}, y_n + \frac{\Delta t}{2} k_2\right) \\
    k_4 &= f(t_n + \Delta t, y_n + \Delta t k_3) 
\end{aligned}
\end{equation}
\begin{equation}
    y_{n+1} = y_n + \frac{\Delta t}{6} (k_1 + 2k_2 + 2k_3 + k_4) + O(\Delta t^5)
\end{equation}
A global truncation error of $O(\Delta t^4)$ permits substantially enlarged temporal steps without triggering geometric instability.

\subsection{Physical Applications: Trajectories and Fields}
Translating abstract numeric frameworks into applied kinematic and electromagnetic simulations requires defining specific governing equations. 

\subsubsection{1D Free-Fall Kinematics}
Consider a spherical object (mass $m$, cross-sectional area $A$) under constant gravitational acceleration $g$, subjected to quadratic aerodynamic drag modeled by the fluid density $\rho$ and drag coefficient $c_d$. The governing kinematic differential equation derived from Newton's second law is:
\begin{equation}
    m \frac{dv}{dt} = -mg + \frac{1}{2} \rho v^2 c_d A
\end{equation}
Discretizing this via the Euler method entails twin iterative updates for position $y$ and velocity $v$:
\begin{equation}
\begin{aligned}
    v_{n+1} &= v_n + \left( -g + \frac{\rho c_d A}{2m} v_n^2 \right) \Delta t \\
    y_{n+1} &= y_n + v_n \Delta t
\end{aligned}
\end{equation}
To pinpoint the exact bounding collision time (e.g., $y=0$) when simulation frames overshoot the spatial boundary condition ($y_{n+1} < 0$), linear interpolation bridges the gap between the pre-collision state ($t_n, y_n$) and the post-collision state ($t_{n+1}, y_{n+1}$):
\begin{equation}
    t_{\text{land}} = t_n + \frac{y_n}{y_n - y_{n+1}} \Delta t
\end{equation}
Compared to purely analytical parabolic solutions applicable only in a vacuum ($y(t) = h - \frac{1}{2}gt^2$), these coupled numerical techniques faithfully simulate realistic atmospheric entry dynamics where strictly analytical inversions fail.

\subsubsection{Electrically Charged Particle Dynamics}
Computing particle motions within intricate multidimensional electric fields entails evaluating the Lorentz force iteratively. The composite E-field vector at position $\vec{r}$ induced by an assembly of discrete point charges $q_i$ localized at $\vec{r}_i$ is defined by Coulomb's superposition:
\begin{equation}
    \vec{E}(\vec{r}) = \frac{1}{4 \pi \epsilon_0} \sum_i q_i \frac{\vec{r} - \vec{r}_i}{|\vec{r} - \vec{r}_i|^3}
\end{equation}
The kinematic acceleration maps directly from the field via $\frac{d\vec{v}}{dt} = \frac{q}{m} \vec{E}(\vec{r})$. Transferred into a numeric solver like RK4, the decoupled vector evolution yields high-fidelity simulations of complex particle scatterings or orbital captures.

\section{Machine Learning in Physics: Regression and Classification}

The statistical principles driving curve fitting extend transparently into modern computational learning theory. The task of generalized predictive modeling bifurcates critically into Regression (mapping to continuous regimes) and Classification (mapping to discrete eigenstates).

\subsection{Linear Regression Foundations}
Regression models extrapolate continuous physical quantities (e.g., thermal conductivities, energies). Analogous to fundamental curve fitting, a multidimensional linear hypothesis operates via an inner product of the parameter weight vector $\vec{\theta}$ and the feature matrix vector $\vec{x}$:
\begin{equation}
    h_{\theta}(\vec{x}) = \vec{\theta}^T \vec{x}
\end{equation}
The optimization explicitly targets the Mean Squared Error (MSE) cost function $J(\vec{\theta})$, structurally identical to the least squares target:
\begin{equation}
    J(\vec{\theta}) = \frac{1}{2N} \sum_{i=1}^{N} \left( h_{\theta}(\vec{x}^{(i)}) - y^{(i)} \right)^2
\end{equation}

\subsection{Classification and Logistic Probability}
Conversely, state classification (e.g., determining whether a phase transition has occurred) outputs discrete binary labels, $y \in \{0, 1\}$. Linear maps violate absolute probability bounds, mandating a non-linear thresholding function. Logistic Regression employs the standard Sigmoid function to map the linear hypothesis onto a strict $[0,1]$ probability distribution:
\begin{equation}
    P(y=1 | \vec{x}; \vec{\theta}) = \sigma(\vec{\theta}^T \vec{x}) = \frac{1}{1 + e^{-\vec{\theta}^T \vec{x}}}
\end{equation}
Because MSE exhibits pathological non-convexity under sigmoid transforms preventing robust gradient descent, classifiers employ the Binary Cross-Entropy (Log-Loss) metric:
\begin{equation}
\begin{split}
    J(\vec{\theta}) = - \frac{1}{N} \sum_{i=1}^{N} \Big[ &y^{(i)} \log(h_{\theta}(\vec{x}^{(i)})) \\
    &+ (1 - y^{(i)}) \log(1 - h_{\theta}(\vec{x}^{(i)})) \Big]
\end{split}
\end{equation}
When differentiated, the gradient of the Log-Loss surprisingly mirrors the linear regression gradient structurally: $\nabla J = \frac{1}{N}\sum \left(h_{\theta}(\vec{x}^{(i)}) - y^{(i)}\right) \vec{x}^{(i)}$, showcasing profound mathematical symmetries connecting continuous and discrete learning architectures.

\section{Root Finding for Nonlinear Equations}

Extracting boundary conditions, energy eigenvalues, or threshold intersections often entails solving arbitrary nonlinear equations implicitly, $f(x)=0$, where algebraic inversion fails. Computational solvers isolate these roots iteratively. While brute-force scanning evaluates the function densely to isolate initial sign-change thresholds across large spaces, advanced algorithms optimize the convergence speed leveraging local topology.

\subsection{Bisection Method}
The Bisection method relies fundamentally on the Intermediate Value Theorem. Given an initial bracket $[a_0, b_0]$ where $f(a_0)f(b_0) < 0$ (indicating a zero-crossing), the algorithm iteratively halves the search space. At each step $n$, the midpoint $c_n = \frac{a_n + b_n}{2}$ is evaluated. The sub-interval containing the root is retained by discarding the boundary that shares the exact same sign as $f(c_n)$. The absolute bound on the error after $n$ iterations is strictly deterministic:
\begin{equation}
    |x_{\text{root}} - c_n| \leq \frac{|b_0 - a_0|}{2^{n+1}}
\end{equation}
While guaranteed to converge (unconditional stability), the linear convergence rate restricts its utility when high-precision roots are required computationally rapidly.

\subsection{Newton-Raphson Method}
To vastly accelerate convergence, the Newton-Raphson method exploits the local gradient (analytical derivative). By applying a first-order Taylor expansion around the current guess $x_n$:
\begin{equation}
    f(x) \approx f(x_n) + f'(x_n)(x - x_n) = 0
\end{equation}
Solving for the linear intercept $(x - x_n)$ yields the recursive update rule:
\begin{equation}
    x_{n+1} = x_n - \frac{f(x_n)}{f'(x_n)}
\end{equation}
Assuming a sufficiently close initial guess and a non-zero derivative at the root ($f'(\alpha) \neq 0$), the Newton-Raphson method boasts exceptional quadratic convergence, effectively doubling the number of correct decimal digits per iteration. However, extracting the analytical derivative $f'(x)$ is often algorithmically impossible for complex numerical physical models.

\subsection{Secant Method}
When the analytical derivative is inaccessible or computationally expensive to extract, the Secant method approximates $f'(x_n)$ utilizing a finite difference scheme constructed from the two most recent iterations:
\begin{equation}
    f'(x_n) \approx \frac{f(x_n) - f(x_{n-1})}{x_n - x_{n-1}}
\end{equation}
Substituting this geometrical approximation back into the Newton iteration maps a secant line between the two latest coordinates:
\begin{equation}
    x_{n+1} = x_n - f(x_n) \frac{x_n - x_{n-1}}{f(x_n) - f(x_{n-1})}
\end{equation}
The Secant method demands two initial conditions ($x_0, x_1$) and avoids bracket containment. Its convergence scaling is superlinear (specifically $\approx 1.618$, the golden ratio), isolating a pristine balance between computational simplicity (no calculus required) and convergence velocity.

\section{Fourier Analysis and Signal Diagnostics}

\subsection{From Fourier Series to Fourier Transform}

The fundamental premise of Fourier analysis is that any physically realizable signal $f(t)$ can be decomposed into a superposition of sinusoidal basis functions. For a periodic signal with period $T$, the Fourier series expansion in complex exponential form is:
\begin{equation}
    f(t) = \sum_{n=-\infty}^{\infty} c_n e^{i n \omega_0 t}, \quad \omega_0 = \frac{2\pi}{T}
\end{equation}
where the discrete spectral coefficients $c_n$ are determined by the projection:
\begin{equation}
    c_n = \frac{1}{T} \int_{-T/2}^{T/2} f(t) e^{-i n \omega_0 t} \, dt
\end{equation}
As $T \to \infty$, the fundamental frequency spacing $\Delta \omega = \omega_0$ becomes infinitesimal $d\omega$, and the discrete summation transitions into a continuous integral. This limit yields the Fourier Transform pair, bridging the gap between periodic steady-states and transient physical phenomena:
\begin{equation}
    f(t) = \frac{1}{2\pi} \int_{-\infty}^{\infty} F(\omega) e^{i\omega t} \, d\omega
\end{equation}
\begin{equation}
    F(\omega) = \int_{-\infty}^{\infty} f(t) e^{-i\omega t} \, dt
\end{equation}

\subsection{Mathematical Formalism and Physical Computation}

In computational physics, the choice of normalization constant and sign convention depends on the specific domain (e.g., Quantum Mechanics vs. Signal Processing). While the $1/2\pi$ factor above is common in physics to preserve the angular frequency $\omega$ relation, the most physically significant result is the conservation of energy, articulated by **Parseval's Theorem**:
\begin{equation}
    \int_{-\infty}^{\infty} |f(t)|^2 \, dt = \frac{1}{2\pi} \int_{-\infty}^{\infty} |F(\omega)|^2 \, d\omega
\end{equation}
This theorem ensures that the total "energy" (or power density) of a signal is invariant under the transformation, a critical requirement for analyzing wave functions $|\psi(x)|^2$ or electrical power spectra. Furthermore, the spread of a packet in time ($\Delta t$) and frequency ($\Delta \omega$) is restricted by the Uncertainty Principle, $\Delta t \Delta \omega \geq 1/2$, which dictates the limits of signal resolution in experimental physics.


\subsection{Discrete and Inverse Fourier Transforms}

Converting the continuous theoretical frameworks into computational algorithms requires discretization. A sampled numeric dataset limits operations to the Discrete Fourier Transform (DFT), where the integral is replaced by a finite summation. This transition is essential for analyzing experimental signals, which are often corrupted by Gaussian noise: $x_{\text{noisy}}(t) = x(t) + \mathcal{N}(0, \sigma^2)$. Defined over $N$ uniform discrete time samples $x_n$, the transformation projects the signal onto orthogonal sinusoidal basis vectors:

\begin{equation}
    X_k = \sum_{n=0}^{N-1} x_n e^{-i \left( \frac{2\pi}{N} \right) k n}
\end{equation}
Conversely, the physical time-domain signal is precisely reconstructed via the Inverse Discrete Fourier Transform (IDFT):
\begin{equation}
    x_n = \frac{1}{N} \sum_{k=0}^{N-1} X_k e^{i \left( \frac{2\pi}{N} \right) k n}
\end{equation}

\subsection{Frequency Resolution and Aliasing}
When employing algorithms like NumPy's FFT, understanding the spectral topology is paramount. The frequency resolution is inherently governed by the sampling interval $\Delta t$ and array size $N$, forming frequency bins $f_k = \frac{k}{N \Delta t}$. Furthermore, the maximum identifiable frequency is rigidly bound by the Nyquist cutoff threshold:
\begin{equation}
    f_{\text{cutoff}} = \frac{1}{2 \Delta t}
\end{equation}
Analyzing a composite waveform, for example, $x(t) = \sin(4\pi t) + 0.5\sin(10\pi t)$ comprising $2\text{ Hz}$ and $5\text{ Hz}$ components, FFT neatly isolates these localized spikes out of random background noise. If an identical signal suffers a time shift (delay), the fundamental power spectrum magnitude $|X_k|$ remains perfectly invariant; the temporal delay manifests exclusively as a linear phase shift across the complex argument.

\section{Conclusion}

The mathematical engines examined—ranging from stochastic sampling to tensor error optimization, discrete calculus, and trigonometric transforms—construct the structural paradigms of modern physical research. These formal tools enable scientists to interrogate computational limits, seamlessly simulating unreplicable high-energy boundary conditions or complex multi-variate statistical fields where analytical mathematics entirely breakdown.



\bibliography{bibliography}

\end{document}